% !TeX program = pdflatex
% =============================================================================
% Arbeitsblatt 3: Strom, Spannung und Widerstand (quantitativ)
% UZH Praktikum I -- Sara Romer Urban, Kantonsschule im Lee, HS2026
% Klasse: 1f (16.09.2026), aus PLAN_Dokumentengenerierung.md.
%
% ENTWURF, NOCH NICHT INS REPO KOPIERT (auf Wunsch von Loris -- er benennt
% zuerst das bestehende Gruppenpuzzle-Dokument um, das bisher denselben
% Dateinamen traegt).
%
% Getrennt von arbeitsblaettern/Aufgabenblatt nach Loris' eigenem Kriterium
% (nicht dem generischen Werkzeug-Kriterium aus institutionen/
% UZH-Praktikum-I/CLAUDE.md): Arbeitsblatt = alles, was gemeinsam im Plenum
% besprochen/hergeleitet wird (reine Theorie, keine Aufgaben); Aufgabenblatt
% = alle Uebungsaufgaben dazu, siehe aufgabenblatt-u-i-r.tex im selben
% Ordner. Diese Aufteilung wurde explizit so besprochen.
%
% Inhalt (nach Absprache): vollstaendige Einheitenherleitung fuer Ampere
% ([I]=1 C/s) und Ohm ([R]=1 V/A, ueber die Definition R:=U/I, an dieser
% Stelle auch das Ohmsche Gesetz selbst), Volt als bereits (aus Batterien/
% Steckdose) bekannte Einheit eingefuehrt statt energetisch hergeleitet
% (U=W/Q waere fuer diese Klasse noch zu frueh). Danach eine explizite
% Schritt-fuer-Schritt-Umformung von U=R*I nach R bzw. nach I (Klasse hat
% laut Vorgabe Muehe mit Umformen), das URI-Dreieck als zusaetzliche
% Kontrollhilfe (nicht als Ersatz fuer die Algebra -- auf Wunsch ggf. auch
% weglassbar). Zum Schluss die Vorsaetze Milli und Kilo (bewusst nur diese
% beiden, keine sehr grossen Vorsaetze wie Mega/Giga). Bewusst NICHT
% enthalten: ein U-I-Diagramm bzw. eine experimentelle/grafische
% Bestimmung von R (auf Wunsch fuer eine spaetere Gelegenheit reserviert).
%
% Bilder: nur 01_Ohmsches-Gesetz-URI-Dreieck-1024x576.jpg wird gebraucht.
% =============================================================================
\documentclass[addpoints,11pt,a4paper]{exam}

\usepackage[T1]{fontenc}
\usepackage[utf8]{inputenc}
\usepackage[ngerman]{babel}
\usepackage{lmodern}
\usepackage[a4paper,margin=2.2cm,headheight=14pt]{geometry}
\usepackage{amsmath}
\usepackage{siunitx}
% Hausstil: zusammengesetzte Einheiten immer als echter Bruch (\frac), nicht
% als Inline-Schraegstrich oder \tfrac -- gilt global fuer jedes \si/\SI.
\sisetup{per-mode=fraction,fraction-command=\frac}
\usepackage{array,booktabs,tabularx,multirow}
\usepackage{enumitem}
\usepackage{xcolor}
\usepackage{colortbl}
\usepackage{tikz}
\usepackage{physics}
\usepackage[most]{tcolorbox}
\usepackage{lastpage}
\usepackage{graphicx}
\usepackage{hyperref}

\usetikzlibrary{arrows.meta,calc,angles,quotes}

\usepackage{setspace}
\usepackage{needspace}
\setlength{\parindent}{0pt}
\setlength{\parskip}{0.6em}
\setstretch{1.08}
\setlist[itemize]{itemsep=0.4em}

% -----------------------------
% Farben (Corporate-Farbschema Physik KFR)
% -----------------------------
\definecolor{titleblue}{HTML}{183A6B}
\definecolor{accentblue}{HTML}{2E6F9E}
\definecolor{lightblue}{HTML}{EAF4FB}
\definecolor{lightgray}{HTML}{F4F6F8}
\definecolor{darkgray}{HTML}{4A4A4A}
\definecolor{accentorange}{HTML}{D9720C}
\definecolor{accentpurple}{HTML}{7B2D8E}
\definecolor{lightorange}{HTML}{FCEEE0}
\definecolor{lightpurple}{HTML}{F3EAF7}

% Links (hyperref): farblich ans Corporate-Farbschema angeglichen.
\hypersetup{
  colorlinks=true,
  urlcolor=accentblue,
  linkcolor=accentblue,
  pdftitle={Arbeitsblatt 3: Strom, Spannung und Widerstand},
  pdfauthor={Loris Keller}
}

% -----------------------------
% Boxen
% -----------------------------
\tcbset{before skip=10pt, after skip=10pt}
\newtcolorbox{titlebox}{
  enhanced,
  colback=titleblue,
  colframe=titleblue,
  boxrule=0pt,
  arc=3mm,
  left=3mm,right=3mm,top=3mm,bottom=3mm
}

% Praktikum I: Lernziele/Einleitung/Theorie werden NICHT farbig gerahmt
% (nur Formel-, Hinweis- und Antwortbox bleiben Boxen) -- siehe
% institutionen/UZH-Praktikum-I/CLAUDE.md, Abschnitt "Boxen-Layout".
\newenvironment{infobox}{%
  \par\addvspace{10pt}\noindent
}{%
  \par\addvspace{10pt}
}

\newenvironment{theoriebox}[1][Theorie]{%
  \par\addvspace{10pt}\noindent\textbf{#1}\par\smallskip\noindent
}{%
  \par\addvspace{10pt}
}

\newtcolorbox{hinweisbox}[1][Hinweis]{
  enhanced,
  breakable,
  colback=lightorange,
  colframe=accentorange,
  title={#1},
  fonttitle=\bfseries,
  boxrule=0.7pt,
  arc=2mm,
  left=5mm,right=5mm,top=2mm,bottom=2mm
}

\newtcolorbox{formelbox}[1][Formel]{
  enhanced,
  breakable,
  colback=lightpurple,
  colframe=accentpurple,
  title={#1},
  fonttitle=\bfseries,
  boxrule=0.9pt,
  arc=2mm,
  left=5mm,right=5mm,top=2mm,bottom=2mm
}

% Bild-Helfer: zentriertes Bild mit spuerbarem Abstand zum umgebenden Text.
% #1 = Breite, #2 = Dateipfad.
\newcommand{\Bild}[2]{%
  \par\vspace{0.6cm}
  \begin{center}
    \includegraphics[width=#1]{#2}
  \end{center}
  \vspace{0.6cm}
}

\pagestyle{headandfoot}
\cfoot{\textcolor{darkgray}{Seite \thepage\ von \pageref{LastPage}}}

\newcommand{\Kopfzeile}[4]{%
  \begin{titlebox}
    {\color{white}\sffamily\bfseries\Large #4}
  \end{titlebox}
  \vspace{0.5em}
  \noindent\textbf{Klasse:}~#1 \hfill \textbf{Datum:}~#2 \hfill \textbf{Thema:}~#3
    \hfill \textbf{Lehrperson:}~Loris Keller
  \vspace{0.8em}
  \lhead{\textcolor{darkgray}{#3}}
  \rhead{\textcolor{darkgray}{#4}}
}

% Werte fuer Kopfzeile zentral setzen
\newcommand{\Klasse}{1f}
\newcommand{\Datum}{16.09.2026}
\newcommand{\Thema}{Ohmsches Gesetz}
\newcommand{\ArbeitsblattTitel}{Arbeitsblatt 3: Strom, Spannung und Widerstand}

\begin{document}

\Kopfzeile{\Klasse}{\Datum}{\Thema}{\ArbeitsblattTitel}

% =============================================================================
% Lernziele
% =============================================================================
\begin{infobox}
\textbf{Lernziele}\\
Am Ende dieser Doppellektion kann ich \dots
\begin{itemize}[leftmargin=1.2cm,itemsep=0.3em]
  \item die Einheiten Ampere, Volt und Ohm nennen und die Einheiten von Strom und Widerstand aus bereits Bekanntem herleiten. (LZ1)
  \item das Ohmsche Gesetz $U=R\cdot I$ anwenden und nach $R$ bzw. nach $I$ umformen. (LZ2)
  \item mit den Vorsätzen Milli und Kilo zwischen verschiedenen Grössenordnungen umrechnen. (LZ3)
\end{itemize}
\end{infobox}

% =============================================================================
% Die Einheiten von Strom, Spannung und Widerstand
% =============================================================================
\needspace{6cm}
\begin{theoriebox}[Die Einheiten von Strom, Spannung und Widerstand]
Bisher haben Sie Strom, Spannung und Widerstand nur qualitativ
kennengelernt. Jetzt lernen Sie, wie man alle drei Grössen tatsächlich
misst und berechnet.
\end{theoriebox}

\needspace{7cm}
\begin{formelbox}[Strom $I$]
Sie kennen bereits $I=\dfrac{\Delta Q}{\Delta t}$ (Ladung pro Zeit). Die
Ladung $Q$ wird in Coulomb $\si{\coulomb}$ gemessen, die Zeit $t$ in
Sekunden $\si{\second}$. Daraus ergibt sich die Einheit des Stroms:
\[
  [I] = 1\ \si{\coulomb\per\second} =: 1\ \text{Ampere} = 1\ \si{\ampere}.
\]
\end{formelbox}

\needspace{5cm}
\begin{formelbox}[Spannung $U$]
Die Spannung wird in Volt $\si{\volt}$ gemessen — diese Einheit kennen
Sie vermutlich bereits von Batterien ($\SI{1.5}{\volt}$, $\SI{9}{\volt}$)
oder von der Steckdose ($\SI{230}{\volt}$).
\end{formelbox}

\needspace{7cm}
\begin{formelbox}[Widerstand $R$]
Der Widerstand ist definiert als das Verhältnis von Spannung zu Strom,
$R:=\dfrac{U}{I}$ (siehe Ohmsches Gesetz unten). Daraus ergibt sich seine
Einheit:
\[
  [R] = 1\ \frac{\si{\volt}}{\si{\ampere}} =: 1\ \text{Ohm} = 1\ \si{\ohm}.
\]
\end{formelbox}

\needspace{8cm}
\begin{theoriebox}[Das Ohmsche Gesetz]
Bei vielen Bauteilen (z.\,B. bei einem handelsüblichen Widerstand) gilt
zwischen Spannung, Widerstand und Strom folgender Zusammenhang, das
\textbf{\emph{Ohmsche Gesetz}}:

\begin{formelbox}[Ohmsches Gesetz]
\[
  U = R\cdot I
\]
$U$: Spannung ($\si{\volt}$); $R$: Widerstand ($\si{\ohm}$); $I$: Strom
($\si{\ampere}$).
\end{formelbox}

Je grösser der Widerstand bei gleichem Strom, desto grösser die Spannung,
die über ihm abfällt — genau wie ein engeres Rohr bei gleichem
Durchfluss einen grösseren Druckunterschied braucht.
\end{theoriebox}

% =============================================================================
% Umformen
% =============================================================================
\needspace{14cm}
\begin{hinweisbox}[Umformen: nach $R$ oder nach $I$ auflösen]
Oft ist nicht $U$, sondern $R$ oder $I$ gesucht. Dazu formen Sie die
Formel $U=R\cdot I$ um, indem Sie auf \emph{beiden} Seiten der Gleichung
dieselbe Rechenoperation ausführen:

\textbf{Nach $R$ auflösen:} Beide Seiten durch $I$ teilen:
\[
  \frac{U}{I} = \frac{R\cdot I}{I} = R \qquad\Longrightarrow\qquad R=\frac{U}{I}.
\]

\textbf{Nach $I$ auflösen:} Beide Seiten durch $R$ teilen:
\[
  \frac{U}{R} = \frac{R\cdot I}{R} = I \qquad\Longrightarrow\qquad I=\frac{U}{R}.
\]

\Bild{7cm}{Bilder/01_Ohmsches-Gesetz-URI-Dreieck-1024x576.jpg}
{\scriptsize Das \glqq URI-Dreieck\grqq{} als Kontrolle: Decken Sie die
gesuchte Grösse ab. Stehen die beiden übrigen Grössen \emph{nebeneinander},
werden sie multipliziert; steht eine \emph{über} der anderen, wird
dividiert. Das ersetzt das Umformen von Hand nicht, hilft aber beim
Nachprüfen.}
\end{hinweisbox}

% =============================================================================
% Grössenordnungen: Vorsätze
% =============================================================================
\needspace{10cm}
\begin{theoriebox}[Grössenordnungen: die Vorsätze Milli und Kilo]
In der Elektrizitätslehre sind die tatsächlich auftretenden Werte oft sehr
klein (Ströme durch kleine Verbraucher, Spannungen in elektronischen
Schaltungen) oder eher gross (Widerstände in Bauteilen). Damit man nicht
ständig mit vielen Nullen rechnen muss, verwendet man
\textbf{\emph{Vorsätze}}, die eine Einheit mit einer Zehnerpotenz
multiplizieren:

\begin{center}
\begin{tabular}{lccc}
\toprule
Vorsatz & Symbol & Faktor & Beispiel \\
\midrule
Milli & m & $10^{-3}$ & $\SI{5}{\milli\ampere} = \SI{0.005}{\ampere}$ \\
(Basiseinheit) & -- & $10^{0}$ & $\SI{1}{\ampere}$, $\SI{1}{\volt}$, $\SI{1}{\ohm}$ \\
Kilo & k & $10^{3}$ & $\SI{2.2}{\kilo\ohm} = \SI{2200}{\ohm}$ \\
\bottomrule
\end{tabular}
\end{center}

Am häufigsten begegnen Ihnen $\si{\milli\ampere}$ (mA, kleine Ströme) und
$\si{\milli\volt}$ (mV, kleine Spannungen) sowie $\si{\kilo\ohm}$
(k$\Omega$, für grössere Widerstände). Um zwischen den Einheiten
umzurechnen, multiplizieren oder dividieren Sie mit dem entsprechenden
Faktor 1000.
\end{theoriebox}

\end{document}
